\section{Conclusion and Future Work}
\label{sec:conclusion}

We have presented GrabDrop, a cross-platform screenshot sharing system
driven by real-time hand gesture recognition. The system demonstrates
that a lightweight TCN model (87K~parameters), combined with
MediaPipe hand landmark detection and a two-stage power-aware
detection pipeline, can achieve practical gesture recognition accuracy
(89.7\%) while remaining deployable on commodity mobile devices through
structured pruning and INT8 quantization (final model size
$\sim$140\,KB).

The project makes several contributions relevant to the course themes
of AI application in practice: (1)~applying cutting-edge temporal
convolutional architectures to a real-world interaction problem;
(2)~demonstrating the full model lifecycle from data collection through
training, optimization, and edge deployment; and (3)~building a
functional cross-platform system with a zero-configuration network
protocol.

\subsection{Limitations}

\begin{itemize}[leftmargin=*]
  \item \textbf{Small test set}: With only 29~test samples, the
    reported accuracy has high variance. A larger, more diverse test
    set would provide more reliable performance estimates.
  \item \textbf{Pruning accuracy loss}: The 30\% structured pruning
    reduced accuracy from 89.7\% to 82.8\%. Knowledge distillation
    from the original model could potentially recover some of this gap.
  \item \textbf{No encryption}: The current network protocol transmits
    screenshots in plaintext. TLS/DTLS encryption should be added for
    production use.
  \item \textbf{Single-hand limitation}: The system detects only one
    hand, precluding two-handed gesture vocabularies.
\end{itemize}

\subsection{Future Directions}

\begin{enumerate}[leftmargin=*]
  \item \textbf{Knowledge distillation}: Train the pruned model under
    supervision of the original model to recover accuracy without
    increasing model size.
  \item \textbf{Expanded gesture vocabulary}: Add gestures such as
    pinch-to-zoom, rotation, or multi-finger taps for richer
    interaction.
  \item \textbf{Federated learning}: Allow on-device personalization
    of the gesture model without sharing raw data.
  \item \textbf{End-to-end encryption}: Implement TLS for TCP
    transfer and DTLS for UDP discovery to secure content in transit.
  \item \textbf{Additional modalities}: Incorporate depth cameras or
    IMU data for more robust gesture recognition in challenging
    environments.
\end{enumerate}
